% Preview source code for paragraph 5

\begin{table}[!h]
\caption{Fi 在不同状态下的表现形式}

\resizebox{\linewidth}{!}{%

\begin{tabular}{ccccccc}
\toprule 
\textbf{功能} & \textbf{角色} & \textbf{状态} & \textbf{涉及类型} & \textbf{触发器} & \textbf{关键字/核心台词} & \textbf{表现形式 ｜ 例子}\tabularnewline
\midrule
\midrule 
\multirow{13}{*}{Fi} & Hero/Parent & 正常 & - & N/A & 本心，真诚，价值 & 忠于自我，内在价值，排斥虚伪\tabularnewline
\cmidrule{2-7}
 & \multirow{2}{*}{Child} & 游戏 & \multirow{2}{*}{\begin{cellvarwidth}[t]
\centering
INTJ\\
ISTJ
\end{cellvarwidth}} & N/A & \multirow{2}{*}{我偏要喜欢} & 傲娇(Tsundere)，深情忠诚，情感纯粹\tabularnewline
\cmidrule{3-3}\cmidrule{5-5}\cmidrule{7-7}
 &  & 幼稚 &  & 压力 &  & 情感任性，无视大局，敏感记仇\tabularnewline
\cmidrule{2-7}
 & \multirow{3}{*}{Anima} & \begin{cellvarwidth}[t]
\centering
正常\\
(Aspiration)
\end{cellvarwidth} & \multirow{3}{*}{\begin{cellvarwidth}[t]
\centering
ENTJ\\
ESTJ
\end{cellvarwidth}} & N/A & ``我也有温柔一面'' & \begin{cellvarwidth}[t]
\centering
隐秘深情，道德洁癖\\
溺爱宠物，对某首老歌流泪，讲义气
\end{cellvarwidth}\tabularnewline
\cmidrule{3-3}\cmidrule{5-7}
 &  & \begin{cellvarwidth}[t]
\centering
压抑\\
(Suppression)
\end{cellvarwidth} &  & 防御性屏蔽 & ``情绪是Te的低效'' & \begin{cellvarwidth}[t]
\centering
情感切除，机器人化\\
``只谈公事''，无视他人痛苦，\\
否认自己有情绪
\end{cellvarwidth}\tabularnewline
\cmidrule{3-3}\cmidrule{5-7}
 &  & \begin{cellvarwidth}[t]
\centering
劫持/抓取\\
(Grip)
\end{cellvarwidth} &  & \begin{cellvarwidth}[t]
\centering
权威丧失，\\
内鬼背叛
\end{cellvarwidth} & ``全世界都对不起我'' & \begin{cellvarwidth}[t]
\centering
陷入自怜自艾(Self-pity)，\\
情绪化爆发，甚至抑郁
\end{cellvarwidth}\tabularnewline
\cmidrule{2-7}
 & \multirow{2}{*}{Nemesis} & \begin{cellvarwidth}[t]
\centering
防御表现\\
(Stubbornness)
\end{cellvarwidth} & \multirow{2}{*}{\begin{cellvarwidth}[t]
\centering
ENFJ\\
ESFJ
\end{cellvarwidth}} & \begin{cellvarwidth}[t]
\centering
主导功能\\
受到挑战
\end{cellvarwidth} & \multirow{2}{*}{``你太自私了''} & \begin{cellvarwidth}[t]
\centering
否定个体：\\
为了群体和谐而压抑个人感受，\\
也看不惯别人的特立独行
\end{cellvarwidth}\tabularnewline
\cmidrule{3-3}\cmidrule{5-5}\cmidrule{7-7}
 &  & \begin{cellvarwidth}[t]
\centering
投射怀疑\\
(Paranoia)
\end{cellvarwidth} &  & \begin{cellvarwidth}[t]
\centering
对他人该功能\\
莫名厌烦
\end{cellvarwidth} &  & \begin{cellvarwidth}[t]
\centering
怀疑人品：\\
认为那些坚持自我价值的人\\
是``破坏团结''的坏人
\end{cellvarwidth}\tabularnewline
\cmidrule{2-7}
 & \multirow{2}{*}{Critic} & 对外打压 & \multirow{2}{*}{INFJ, ISFJ} & 防御机制 & \multirow{2}{*}{批判道德: 你真坏} & \begin{cellvarwidth}[t]
\centering
道德审判，\\
指责自私虚伪冷血人品
\end{cellvarwidth}\tabularnewline
\cmidrule{3-3}\cmidrule{5-5}\cmidrule{7-7}
 &  & 对内自毁 &  & 防御机制 &  & \begin{cellvarwidth}[t]
\centering
自我羞耻，\\
觉得自己是伪君子
\end{cellvarwidth}\tabularnewline
\cmidrule{2-7}
 & Trickster & 忽视/贬低 & \begin{cellvarwidth}[t]
\centering
ENTP\\
ESTP
\end{cellvarwidth} & 无意识的黑洞 & \begin{cellvarwidth}[t]
\centering
``这也太矫情了，\\
客观一点行不行？''
\end{cellvarwidth} & \begin{cellvarwidth}[t]
\centering
情感逻辑化，\\
不知道自己真实感受，\\
无意间把别人的价值观踩得粉碎
\end{cellvarwidth}\tabularnewline
\cmidrule{2-7}
 & \multirow{2}{*}{Demon} & 毁灭一切 & \multirow{2}{*}{\begin{cellvarwidth}[t]
\centering
ENTP\\
ESTP
\end{cellvarwidth}} & \begin{cellvarwidth}[t]
\centering
彻底崩溃\\
严重侵犯\\
极度绝望
\end{cellvarwidth} & \begin{cellvarwidth}[t]
\centering
道德虚无\\
 Nihilism
\end{cellvarwidth} & 反社会倾向，情感报复\tabularnewline
\cmidrule{3-3}\cmidrule{5-7}
 &  & 整合 &  & N/A & \begin{cellvarwidth}[t]
\centering
慈悲圣徒\\
 Universal Love
\end{cellvarwidth} & \begin{cellvarwidth}[t]
\centering
浪子回头般的真诚，\\
对他人的苦难产生深刻共情，\\
为了信仰不惜一切
\end{cellvarwidth}\tabularnewline
\bottomrule
\end{tabular}

}
\end{table}
